\section{A physical same-ultra saving after orbit slicing}
\label{sec:tpc35-physical-ultra}

The preceding sections diagnose abstract routes.  We now return to the
actual multiplier \eqref{eq:tpc35-actual-multiplier}, coefficient vectors,
and ultra increments.  Orbit slicing removes one time sum from the
shared-divisor estimate of TPC-34 and yields a genuine physical saving.

\subsection{Opening the divisor copies}

For the left slice \eqref{eq:tpc35-left-slice}, define
\begin{equation}
 Z_{\gamma,j,u}^L
 :=a_R(u)
 \sum_{\substack{\alpha\\u\mid N_\alpha(j)}}
 \gamma_\alpha^{(1)}
 \mathfrak A_{\alpha,\gamma}^{\rm act}(j),
 \qquad T<u\le U_0.
 \label{eq:tpc35-left-u-slice}
\end{equation}
Then
\begin{equation}
 Y_{\gamma,j}^L=\sum_uZ_{\gamma,j,u}^L.
 \label{eq:tpc35-Y-as-u-sum}
\end{equation}
The diagonal in the two opened ultra-divisor copies is
\begin{equation}
 \cV_{L,u=}:=\sum_{\gamma,j}\sum_{T<u\le U_0}
 |Z_{\gamma,j,u}^L|^2.
 \label{eq:tpc35-left-same-u}
\end{equation}
Define \(Z_{\alpha,j,u}^R\) and \(\cV_{R,u=}\) by exchanging the two row
systems.  These quantities are nonnegative.  They are not, however,
lower or upper bounds for \(V_L,V_R\), because the distinct-\(u\) cross
terms in \eqref{eq:tpc35-Y-as-u-sum} may have either sign.

For a termwise absolute comparison with TPC-34, also define
\begin{equation}
 \widetilde{\cV}_{L,u=}
 :=\sum_{\gamma,j}\sum_{T<u\le U_0}|a_R(u)|^2
 \left(
 \sum_{\substack{\alpha\\u\mid N_\alpha(j)}}
 \left|
 \gamma_\alpha^{(1)}
 \mathfrak A_{\alpha,\gamma}^{\rm act}(j)
 \right|
 \right)^2,
 \label{eq:tpc35-left-absolute-same-u}
\end{equation}
and symmetrically \(\widetilde{\cV}_{R,u=}\).  The triangle inequality
gives
\begin{equation}
 \cV_{L,u=}\le\widetilde{\cV}_{L,u=},
 \qquad
 \cV_{R,u=}\le\widetilde{\cV}_{R,u=}.
 \label{eq:tpc35-same-u-majorization}
\end{equation}

\subsection{Residue occupancy at one time}

\begin{lemma}[Fixed-\((u,j)\) row occupancy]
\label{lem:tpc35-fixed-uj-occupancy}
On the physical support,
\begin{equation}
 \#\{\alpha:u\mid N_\alpha(j)\}
 \ll_{\mathscr D}1+\frac Qu
 \le 1+\frac QT.
 \label{eq:tpc35-fixed-uj-occupancy}
\end{equation}
\end{lemma}

\begin{proof}
Primitivity \eqref{eq:tpc35-divisor-primitivity} gives \((u,j)=1\).
Hence
\[
 u\mid m_\alpha j+h
 \quad\Longleftrightarrow\quad
 m_\alpha\equiv-hj^{-1}\pmod u.
\]
The distinct row labels \(m_\alpha\) lie in an interval of length
\(O_{\mathscr D}(Q)\), so one residue class modulo \(u\) contains
\(O_{\mathscr D}(1+Q/u)\) of them.  Since \(u>T\), the second bound follows.
\end{proof}

\begin{theorem}[\(T\)-saving bound for the opened shared-ultra block]
\label{thm:tpc35-shared-ultra}
On the physical packet,
\begin{equation}
 \boxed{
 \cV_{L,u=}+\cV_{R,u=}
 \le\widetilde{\cV}_{L,u=}+\widetilde{\cV}_{R,u=}
 \ll_{\eps,h,\mathscr D}
 X^\eps Q^2J\left(1+\frac QT\right).}
 \label{eq:tpc35-shared-ultra-bound}
\end{equation}
At the endpoint \eqref{eq:tpc35-high-beta-scales},
\begin{equation}
 \boxed{
 \cV_{L,u=}+\cV_{R,u=}
 \le\widetilde{\cV}_{L,u=}+\widetilde{\cV}_{R,u=}
 \ll X^{o(1)}\frac{Q^3J}{T}.}
 \label{eq:tpc35-shared-ultra-endpoint}
\end{equation}
\end{theorem}

\begin{proof}
For fixed \(\gamma,j,u\), Cauchy--Schwarz and
\cref{lem:tpc35-fixed-uj-occupancy} give
\begin{align*}
 &|a_R(u)|^2
 \left(
 \sum_{\substack{\alpha\\u\mid N_\alpha(j)}}
 \left|\gamma_\alpha^{(1)}
 \mathfrak A_{\alpha,\gamma}^{\rm act}(j)\right|
 \right)^2\\
 &\quad\le |a_R(u)|^2
 \#\{\alpha:u\mid N_\alpha(j)\}
 \sum_{\substack{\alpha\\u\mid N_\alpha(j)}}
 |\gamma_\alpha^{(1)}|^2
 |\mathfrak A_{\alpha,\gamma}^{\rm act}(j)|^2\\
 &\ll X^\eps\left(1+\frac QT\right)
 \sum_{\substack{\alpha\\u\mid N_\alpha(j)}}
 |\gamma_\alpha^{(1)}|^2.
\end{align*}
Here logarithmic factors in \(a_R\) and the bounded multiplier are
absorbed into \(X^\eps\).  Sum first over \(u\).  For each fixed
\((\alpha,j)\), the number of divisors of \(N_\alpha(j)\asymp X\) is
\(O_\eps(X^\eps)\).  Then sum over \(O(Q)\) opposite rows \(\gamma\), over
\(O(J)\) orbit times, and use \eqref{eq:tpc35-row-l2}.  This yields
\[
 \widetilde{\cV}_{L,u=}
 \ll X^\eps\left(1+\frac QT\right)QJ\cdot Q.
\]
The right side is identical.  Combine with
\eqref{eq:tpc35-same-u-majorization}.  Since \(Q/T\to\infty\) at the
endpoint, \eqref{eq:tpc35-shared-ultra-endpoint} follows.
\end{proof}

The absolute orbit--sliced baseline is \(Q^3J\), so
\eqref{eq:tpc35-shared-ultra-endpoint} saves the full factor \(T\).  It is
also one factor \(J\) below the two--time shared-ultra scale
\(Q^3J^2/T\) proved in TPC-34.  Nevertheless
\begin{equation}
 \frac{Q^3J/T}{Q^3/J}=\frac{J^2}{T}
 =X^{279/1000+o(1)},
 \label{eq:tpc35-shared-ultra-gap}
\end{equation}
so the block does not close the orbit gate.

\subsection{Terminal atoms on the same-\texorpdfstring{\(u\)}{u} diagonal}

On the endpoint packet, the terminal divisor
\begin{equation}
 u=N_\alpha(j)=m_\alpha j+h
 \label{eq:tpc35-terminal-divisor}
\end{equation}
lies in the ultra range.  At one fixed orbit time, equal terminal divisors
have no off--diagonal row solutions.

\begin{proposition}[Terminal equality is the row diagonal]
\label{prop:tpc35-terminal-equality}
If \(j>0\) and
\begin{equation}
 N_{\alpha_1}(j)=N_{\alpha_2}(j),
 \label{eq:tpc35-terminal-equality}
\end{equation}
then \(m_{\alpha_1}=m_{\alpha_2}\), hence
\(\alpha_1=\alpha_2\) under the row-label convention.
\end{proposition}

\begin{proof}
Subtracting the two equalities gives
\((m_{\alpha_1}-m_{\alpha_2})j=0\).  Since \(j>0\), the row labels agree.
\end{proof}

Thus the terminal--terminal contribution inside the same-\(u\) block is
already contained in the proved Gram--copy diagonal.  This does not remove
the terminal obstruction identified in TPC-34: two distinct terminal
divisors occur in the cross-\(u\) part of
\(
 |\sum_uZ_{\gamma,j,u}^L|^2
\), where the fused signs produce an averaged four--M\"obius row
correlation.  Nor does the triangle inequality between terminal and
proper-divisor layers give a converse bound, because those layers may
cancel before squaring.  The theorem here proves a \(T\)-saving bound for
one opened positive block and locates the remaining signed interaction
more precisely; it does not break the parity boundary.
